%==============================================================================
%  PlantMetWiki -- Scientific Data ARTICLE
%  Submitted to the Collection: "10 years of the FAIR principles"
%  https://www.nature.com/collections/eedeeaaihi   (deadline 14 November 2026)
%
%  MANUSCRIPT TYPE: Article (the collection accepts Articles ONLY -- not
%  Data Descriptors). Scientific Data Articles "cover data policy, repositories,
%  standards, ontologies, workflows, or any topic relating to the mechanics of
%  data sharing", and unlike Data Descriptors they DO carry Results and
%  Discussion sections. The structure below therefore follows the preprint.
%
%  COMPILER: works with pdfLaTeX, XeLaTeX or LuaLaTeX (see iftex block below).
%  In Overleaf: Menu -> Compiler. LuaLaTeX gives the best Unicode handling for
%  the Greek letters and chemical formulae used throughout.
%==============================================================================
\documentclass[11pt,a4paper]{article}

\usepackage{iftex}
\ifPDFTeX
  \usepackage[utf8]{inputenc}
  \usepackage[T1]{fontenc}
  \usepackage{textcomp}
  \usepackage{newunicodechar}
  % Unicode characters used in this manuscript, mapped for pdfLaTeX.
  \newunicodechar{α}{\ensuremath{\alpha}}
  \newunicodechar{β}{\ensuremath{\beta}}
  \newunicodechar{γ}{\ensuremath{\gamma}}
  \newunicodechar{δ}{\ensuremath{\delta}}
  \newunicodechar{ε}{\ensuremath{\epsilon}}
  \newunicodechar{→}{\ensuremath{\rightarrow}}
  \newunicodechar{×}{\ensuremath{\times}}
  \newunicodechar{≥}{\ensuremath{\geq}}
  \newunicodechar{≤}{\ensuremath{\leq}}
  \newunicodechar{~}{\ensuremath{\sim}}
  \newunicodechar{₂}{\textsubscript{2}}
  \newunicodechar{₃}{\textsubscript{3}}
  \newunicodechar{₀}{\textsubscript{0}}
  \newunicodechar{₆}{\textsubscript{6}}
  \newunicodechar{⁺}{\textsuperscript{+}}
  \newunicodechar{′}{'}
  \newunicodechar{″}{''}
\else
  \usepackage{fontspec}
\fi

\usepackage[margin=2.5cm]{geometry}
\usepackage{graphicx}
\usepackage{booktabs}
\usepackage{longtable}
\usepackage{array}
\usepackage{xcolor}
\usepackage{lineno}
\usepackage[hidelinks]{hyperref}
\usepackage{setspace}

\linenumbers
\onehalfspacing

% Author-facing notes. Delete all \note occurrences before submission.
\newcommand{\note}[1]{\textcolor{red}{\textbf{[NOTE: #1]}}}

\setlength{\parindent}{0pt}
\setlength{\parskip}{0.6em}

\title{\vspace{-2em}\bfseries PlantMetWiki: a FAIR knowledge graph for plant metabolic pathway cross-species representation and integration}
\date{}

\begin{document}

%==============================================================================
% FRONT MATTER
%==============================================================================
\maketitle
\vspace{-3em}

\begin{center}
Elena Del Pup\textsuperscript{1,*},
Max Muller\textsuperscript{1},
Marvin Martens\textsuperscript{2},
Egon L. Willighagen\textsuperscript{2},
Marnix H. Medema\textsuperscript{1,3},
Denise Slenter\textsuperscript{2,4,5,*},
Justin J.J. van der Hooft\textsuperscript{1,6,*}
\end{center}

\vspace{0.5em}
\noindent\footnotesize
\textsuperscript{1}Bioinformatics Group, Wageningen University, Wageningen, the Netherlands.\\
\textsuperscript{2}Department of Translational Genomics, NUTRIM, Maastricht University, Maastricht, the Netherlands.\\
\textsuperscript{3}Institute of Biology, Leiden University, Leiden, the Netherlands.\\
\textsuperscript{4}Department of Advanced Computing Sciences, DACS, Maastricht University, Maastricht, the Netherlands.\\
\textsuperscript{5}Institute of Data Science, IDS, Maastricht University, the Netherlands.\\
\textsuperscript{6}Department of Biochemistry, University of Johannesburg, Johannesburg 2006, South Africa.

\vspace{0.5em}
\noindent\textsuperscript{*}Correspondence may be addressed to: EDP. Email: \href{mailto:elena.delpup@wur.nl}{elena.delpup@wur.nl}.
DS. Email: \href{mailto:denise.slenter@maastrichtuniversity.nl}{denise.slenter@maastrichtuniversity.nl}.
JJJvdH. Email: \href{mailto:justin.vanderhooft@wur.nl}{justin.vanderhooft@wur.nl}
\normalsize

\vspace{1em}
\noindent\textbf{Keywords}: knowledge graphs, FAIR, linked data, plant specialized metabolism, multi-omics

\vspace{1.5em}

%==============================================================================
% ABSTRACT -- restored to the full preprint text (216 words).
% NOTE: the 170-word cap is documented for Data Descriptors. Confirm the limit
% that applies to the Article manuscript type before submission.
%==============================================================================
\section*{Abstract}

Plants produce a vast diversity of specialized metabolites with extensive potential ecological,
agro-industrial, and pharmaceutical applications. Discovery of novel plant natural products relies on
combining multi-omics evidence with biochemical transformations. However, pathway-level annotations
are fragmented across individual species, databases, and publications, limiting comparative
cross-species pathway analyses and systematic generation of hypotheses. To support integration and
reuse of plant metabolic knowledge, we developed PlantMetWiki, a FAIR Linked Open Data semantically
enriched knowledge graph built on infrastructure adapted from WikiPathways. Our approach extends on
the highly curated pathway information from Plant Metabolic Network with crosslinks to biosynthetic
gene clusters resources (MIBiG and plantiSMASH), metabolite annotations, and cross-species modelling.
This way, our resource captures pathway genomic context, increases metabolomics data interoperability
via federated queries, and supports cross-species analysis to identify annotation gaps. PlantMetWiki
represents 1,162 plant metabolic pathways as Resource Description Framework (RDF) graphs, preserving
pathway structure, literature provenance, annotations, and taxonomic information from PlantCyc 17.0.
PlantMetWiki is distributed through a public SPARQL endpoint with open-source reproducible data
transformation and validation workflows. By modelling pathways as a multispecies graph, PlantMetWiki
enables comparative analyses across taxa, integration with external chemical knowledge resources
through federated queries, and identification of metabolic, genomic, and chemical annotation gaps. As
a result, PlantMetWiki provides a foundation for FAIR reuse and integration of plant pathway
knowledge.

%==============================================================================
% INTRODUCTION -- restored in full from the preprint (no shortening).
%==============================================================================
\section*{Introduction}

Plants encompass hundreds of thousands of species that have evolved diverse metabolic repertoires to
adapt to distinct ecological niches\cite{huang2023}. As a result, different plant lineages produce a
vast number of unique combinations of endogenous specialized metabolites, ranging from alkaloids and
terpenoids to phenylpropanoids and glycosides, many of which have important ecological, agricultural,
industrial, and pharmaceutical applications\cite{domingo2023,owen2026,greco2026}. As evolution of
plant chemical diversity largely relies on shared metabolic networks and common
branchpoints\cite{ji2024}, mapping biosynthetic pathways across species remains a central challenge
in plant biology and natural product research.

While the increasing number of experimentally characterized plant metabolic pathways provides
valuable insights into enzymes, metabolites, and regulatory mechanisms underlying specialized
metabolism, biosynthetic knowledge remains unevenly distributed across species, pathways, and
specialized domain-specific resources. Model organisms and economically important crops are
extensively annotated; however, many species contain only partial pathway information, and evidence
for individual pathway steps is often scattered across multiple taxa. Existing databases capture
complementary yet disconnected aspects of metabolism, including enzymatic reactions, metabolite
structures, or genomic context, but are rarely integrated. While these resources are individually
valuable, they are rarely represented within a common framework to allow pathway, genomic, taxonomic,
and chemical information to be explored simultaneously. For instance, major natural product chemical
structure databases such as LOTUS\cite{rutz2022} and COCONUT\cite{sorokina2021} contain structures
and associated metadata but lack explicit links to metabolic pathways, thereby hindering
cross-species analyses supporting enrichment of multi-omics datasets with metabolite information or
annotation of cross-species metabolic networks\cite{mildau2025,singh2022}. Integrating this
fragmented knowledge is essential for comparative pathway analysis, biosynthetic hypothesis
generation, and the discovery of novel plant natural products. Comparative analyses across species
are also essential to translate knowledge from model plant species to less-studied taxa and
crops\cite{uauy2025}, as well as to understand the evolutionary processes that shape specialized
metabolism. Pathways may diversify through lineage-specific innovations, recruit alternative enzymes
to catalyze equivalent reactions, or arise independently through convergent evolutionary strategies
that lead to the production of similar metabolites in distantly related species\cite{scossa2026}. The
relationships between species are captured by taxonomic information, where NCBI
Taxonomy\cite{schoch2020} acts as a reference ontology. Despite the importance of these comparative
perspectives, most pathway resources and their semantic representation remain centered on
species-specific representations, limiting the ability to systematically query shared pathway
components and integrate evidence across biological and omics layers. Cross-species pathway models
are a promising framework to connect these different biological entities together in a harmonized
overview with many options for data visualization, analysis, and interpretation.

The Plant Metabolic Network (PMN)\cite{hawkins2025} provides the extensively curated pan-plant
reference database PlantCyc, which includes machine-readable annotations of enzymes and metabolites
across species-specific biochemical conversions across hundreds of plant taxa. PMN adopts the
BioPAX\cite{demir2010} standard to facilitate pathway data exchange and interoperability and the
Pathway Tools ecosystem\cite{karp2021}. This standardized format, however, presents limitations for
integrative analyses, since the flat tabular representations simplify access at the cost of losing
the rich relational structure of metabolic networks. Consequently, complex relationships such as
hierarchical pathway organization, many-to-many associations between reactions and metabolites, and
detailed provenance information are difficult to represent and query. Furthermore, this format cannot
accommodate integration with other annotations, including biosynthetic gene clusters, predicted
pathways, individual multi-omics datasets, and external knowledge graphs, restricting large-scale
comparative and cross-domain analyses.

The Graphical Pathway Markup Language (GPML)\cite{kutmon2015} provides a pathway modeling standard
that is more flexible in nature and can be extended with new annotations. This format can be
transferred to the Resource Description Framework (RDF), providing an open, community-curated
platform for pathway representation that adopts Semantic Web technologies using vocabularies enabling
data sharing and reuse, persistent identifiers, and SPARQL for querying\cite{waagmeester2016}. This
graph-based approach enables pathways to be represented as interconnected networks, preserving
complex biological relationships while supporting integration with external resources through Linked
Open Data (LOD) for machine-readable interlinked data. The infrastructure developed by
WikiPathways\cite{agrawal2024} facilitates visualization, editing, and analysis of pathway networks,
and supports integration with multi-omics tools and network analysis platforms such as
Cytoscape\cite{shannon2003}, as well as federated queries across databases including
Wikidata\cite{waagmeester2020}. In plants, the use of WikiPathways was piloted to analyze rice and
Arabidopsis seed development networks and gene regulatory networks\cite{hanumappa2013}. However, the
current GPML pathway model used in WikiPathways is constrained to single-species pathway
representations, thereby limiting integration approaches that encompass multiple species and the
genome-scale annotations available in PMN.

Emerging evidence of the physical clustering of plant biosynthetic genes in the chromosome could be
used to annotate gaps in plant metabolic knowledge by integrating these genomic regions called
biosynthetic gene clusters (BGCs) with pathway models. MIBiG\cite{zdouc2025} is a repository of
experimentally validated BGCs which relies on a minimum information standard for their annotation,
continuously powered by collaborative community annotations, and contains 43 validated plant BGCs in
version 4.0. However, MIBiG lacks representation of partially clustered plant pathways, which are
abundant in plants across taxa\cite{itkin2013}. Additionally, the MIBiG schema does not capture
multispecies information for BGCs that occur in several species or that are similar to each other,
despite many key studies for their discovery and characterization rely on comparative
approaches\cite{franke2019,li2021}. This limits follow-up comparative studies on the relationships
between clustered and unclustered biosynthesis. The plantiSMASH\cite{delpup2026} database captures
more than 30,000 predicted plant BGCs across more than 400 species, detected by mining plant genomes
for enzyme domains associated with known biosynthetic ``rules''. The database maps
plantiSMASH-detected BGCs to validated ones from MIBiG and to other similar putative BGCs across
species, facilitating cross-species analyses. However, it does not provide links of
plantiSMASH-detected BGCs with wider metabolic networks.

To overcome the above-mentioned gaps, here we present Plant Metabolic Pathways Wiki (PlantMetWiki), a
LOD knowledge graph for representing and exploring plant metabolic pathways across species.
PlantMetWiki combines curated pathway knowledge from PlantCyc with the semantic infrastructure of
WikiPathways\cite{agrawal2024}, creating a unified graph-based representation of plant metabolism
that preserves pathway structure while enabling interoperability with external resources. The
resource integrates experimentally characterized pathway annotations from PlantCyc 17.0, taxonomic
information from NCBI, metabolite identifiers through mappings via
BridgeDb\cite{slenter2018,vaniersel2010}, crosslinks to biosynthetic gene cluster information from
MIBiG 4.0\cite{zdouc2025} and predicted clusters from plantiSMASH 2.0\cite{delpup2026}, and natural
product chemistry through federated queries to Wikidata.

PlantMetWiki enables users to (1) compare pathways across species, leveraging shared reactions,
enzymes, substrates, and products, to identify annotation gaps, (2) explore connections between
curated pathways and their genomic context, including predictions from plantiSMASH genome mining, and
(3) calculate completeness between complementary data resources in the natural products space with
federated queries. In this work, we describe the generation of PlantMetWiki, its FAIR (Findable,
Accessible, Interoperable, and Reusable) data infrastructure, and example applications demonstrating
cross-species pathway comparison and integration with plant natural product knowledge. By
representing pathways as a multispecies semantic graph with (gene, metabolite, and taxa) identifier
mapping and RDF-based data models available for querying, PlantMetWiki connects pathway information
to complementary biological, taxonomic, and chemical knowledge available across the LOD ecosystem and
supports query-driven analyses. Additionally, our resource provides an expandable semantic model and
a reproducible framework that can be transferred to other kingdoms of life using multispecies models
for pathway knowledge gap discovery, such as microbes and fungi.

%==============================================================================
% RESULTS
%==============================================================================
\section*{Results}

PlantMetWiki is accessible through an open-access SNORQL User Interface (UI)
(\url{https://plantmetwiki.bioinformatics.nl}) and a dedicated SPARQL endpoint
(\url{https://plantmetwiki.bioinformatics.nl/sparql}), including diverse customizable example
queries. A dedicated tutorial includes relevant documentation of the UI, endpoint, and example
queries, supporting integration with user-provided data, including multi-omics data, phenotypes,
federated queries to Wikidata, and BGCs (from MIBiG and plantiSMASH). Results can be exported in
standard formats such as CSV, JSON, and XML to be further analyzed using network tools such as
Cytoscape and combined with multi-omics visualizations and integrative representations such as
metabolic networks\cite{mildau2025} or paired omics approaches\cite{wolters2024}.

Querying the semantically modelled pathway information in PlantMetWiki supports identification of
curation gaps and prioritization of candidate annotations through alignment of cross-species pathway
information and independent knowledge resources, crosslinks in the knowledge graph (to MIBiG, NCBI,
plantiSMASH), and federated queries. PlantMetWiki supports the FAIR principles by adopting RDF-based
semantic representations, including standard ontologies, machine-readable metadata, and public data
access mechanisms, described in more detail in Table~\ref{tab:fair}.

\begin{longtable}{p{0.44\textwidth} p{0.50\textwidth}}
\caption{Adaptation of the FAIR principles in the PlantMetWiki resource.}
\label{tab:fair}\\
\toprule
\textbf{FAIR principle} & \textbf{Implementation in PlantMetWiki} \\
\midrule
\endfirsthead
\toprule
\textbf{FAIR principle} & \textbf{Implementation in PlantMetWiki} \\
\midrule
\endhead
\bottomrule
\endfoot
F1: (Meta)data are assigned globally unique and persistent identifiers &
(1) Uniform Resource Identifiers (URIs) resource identifiers for all data elements;
(2) Resolvable Uniform Resource Locators (URLs) for publications and metabolites for all species;
(3) Identifier mapping of genes and proteins with resolvable URLs in \textit{Arabidopsis thaliana} \\
\addlinespace
F2: Data are described with rich metadata (defined by R1 below) &
RDF model is supported with VoID header file, including license and provenance information. \\
\addlinespace
F3: Metadata clearly and explicitly include the identifier of the data they describe &
Identifier, version, and location of releases from PlantCyc, plantiSMASH, MIBiG, NCBI ontology, and
BridgeDb metabolite mapping files are available in the VoID header file. \\
\addlinespace
F4: (Meta)data are registered or indexed in a searchable resource & SPARQL endpoint \\
\addlinespace
A1: (Meta)data are retrievable by their identifier using a standardized communications protocol.
A1.1 The protocol is open, free, and universally implementable.
A1.2 The protocol allows for an authentication and authorization procedure, where necessary &
(1) Standardized access via a SPARQL endpoint;
(2) SPARQL documentation is open, free, and universally implementable, enhanced with a dedicated
tutorial;
(3) Authentication and authorization can be implemented, however not necessary due to the nature of
our data \\
\addlinespace
A2: Metadata are accessible, even when the data are no longer available &
(1) All data accessible via the SPARQL endpoint are archived on Zenodo;
(2) The software to create the data is available on GitHub \\
\addlinespace
I1: (Meta)data use a formal, accessible, shared, and broadly applicable language for knowledge
representation &
(1) Framework is based on 10+ years RDF model from WikiPathways;
(2) VoID header file follows the Dublin Core Schema \\
\addlinespace
I2: (Meta)data use vocabularies that follow FAIR principles &
Included vocabularies are based on OWL representations available on BioPortal (WikiPathways) and OBO
Foundry (NCBITaxon ontology) \\
\addlinespace
I3: (Meta)data include qualified references to other (meta)data &
Interactions between biological entities are described using the new GPML2021 schema, including
detailed biological interactions compatible with cross-references \\
\addlinespace
R1: (Meta)data are richly described with a plurality of accurate and relevant attributes.
R1.1 (Meta)data are released with a clear and accessible data usage license.
R1.2 (Meta)data are associated with detailed provenance.
R1.3 (Meta)data meet domain-relevant community standards &
(1) License and provenance are recorded in the VoID header file for automated retrieval;
(2) License and provenance of data and software are included in the SNORQL UI banner, to support user
awareness;
(3) GPML and BridgeDb are recognized resources in FAIRSharing, on which the PlantMetWiki tool is
built \\
\end{longtable}

\subsection*{Building a cross-species semantic model and expandable RDF framework for plant specialized metabolism}

The PlantMetWiki open-source web interface is the result of six co-designed and jointly operating
repositories hosted in the GitHub Pathway-LOD community (see Code availability statement) that
capture all steps from transforming data from PlantCyc in their native BioPAX standard to GPML2021
and RDF, metabolite mappings via BridgeDb, establishing BGC crosslinks, creating a metadata schema,
versioning, automated SNORQL UI hosting, SPARQL example querying, documentation, and tutorials
(Figure~\ref{fig:pipeline}). Data, metadata and software versioning and archiving on Zenodo are
automated steps, with dedicated workflows available on GitHub (see Code availability statement).
Pathway and reaction data from PlantCyc 17.0 are modelled into 1,162 multispecies pathway files and
1,316 individual enzymatic reaction files in the GPML2021 format, then converted to one RDF Turtle
file using the WikiPathways gpml2rdf tool (v4.0.4), providing harmonized representations of all
pathways and associated reactions. The pathways captured by both file types are represented as
\texttt{wp:Pathway} objects and together constitute 2,478 pathway instances in the resulting core
PlantMetWiki knowledge graph. Pathways contain different entity types (DataNodes), modelled following
the WikiPathways vocabulary (\url{http://vocabularies.wikipathways.org/wp#}), representing 2,758
genes (\texttt{wp:GeneProduct}), 3,750 enzymes (\texttt{wp:Protein}), 70 protein complexes
(\texttt{wp:Complex}), 4,577 metabolites (\texttt{wp:Metabolite}), and 54,800 Interactions
(Fig.~S1). Interactions encode directed biological events between DataNodes as
\texttt{wp:target}/\texttt{wp:source} edges (Figure~\ref{fig:model}), and are split into 15 subtypes,
summarized in Fig.~S2. Conversions (\texttt{wp:Conversion}) representing biochemical reactions
involving metabolite substrates and products encode 36.4\% (19,927) of the total Interactions,
followed by Catalysis (\texttt{wp:Catalysis}), linking \texttt{wp:GeneProduct} or \texttt{wp:Protein}
to the reaction these entities catalyze, with 22.6\% (12,359),
\texttt{wp:TranscriptionTranslation} (links genes to the protein these encode) 15.1\% (8,280),
\texttt{wp:Inhibition}/\texttt{wp:Stimulation} (negative or positive regulatory effects) 6.4\%
(3,486) and 1.6\% (865), respectively, \texttt{wp:Binding}/\texttt{wp:ComplexBinding} (physical
interaction or complex formation) 0.1\% (70). Some PlantCyc pathway edges, particularly transport and
regulatory links, are translated to \texttt{wp:DirectedInteraction} (44,934, 82\% total Interaction
coverage) as a fallback type when no more specific class applies. The cumulative coverage of entities
per pathway shows that a few large pathways dominate the resource, with genes and proteins cumulative
curves across the entire database rising steeply (Fig.~S3). Metabolites do not follow such a steep
trend, since their nodes are often shared across pathways. The largest pathways in terms of genes,
metabolites, and species represented from PlantCyc include super pathways of flavones and derivatives
(PWY-6266), super pathway of anaerobic sucrose degradation (PWY-7345), and quercetin glycoside
biosynthesis (PWY-5321) (Fig.~S4--5).

\note{Insert preprint Results lines 285--433 here (taxonomy modelling: ORG-ID to NCBI mapping,
validation outcomes, the taxonomy graph, BGC crosslinks, and the named-graph organisation of the
triplestore). Copy from your source manuscript rather than retyping from the PDF.}

\note{CONSISTENCY -- please reconcile before submission. Across the preprint these numbers disagree:
(i) total triples given as ``$\sim$28.2 million'' (Fig.~1 caption) while the seven named graphs as
listed sum to 6,500,299, and the live endpoint currently returns 31,818,800;
(ii) the NCBITaxon graph is given as 21.7M triples in the Fig.~1 caption but 17,707 in the Fig.~1
graph box and in the Methods (after ROBOT MIREOT subsetting);
(iii) BGC--pathway crosslinks given as 199 in the Fig.~1 box but 225 in the Results;
(iv) taxa given as 439 (PlantCyc source), 424 (taxonomy graph), and ``all 439'' in the Fig.~1 caption
against ``424 seed taxa'' in the Fig.~1 box;
(v) the taxonomy graph is 30,176 triples in the Results and 30,178 in Fig.~1.}

\subsection*{Multispecies pathway modelling enables comparative cross-species pathway analysis and annotation prioritization}

A central feature of PlantMetWiki is the representation of pathways as multispecies semantic graphs
in which genes, enzymes, reactions, metabolites, and taxonomic annotations can be queried together.
This representation enables comparison of pathway coverage across species (e.g., via the previously
mentioned catalysis-chain assignment) and supports the identification of annotation gaps that may not
be apparent in species-specific pathway views. For each pathway, we classify missing species-reaction
annotation as an \emph{inferable gap} if at least one other species in the same pathway has a direct
enzyme annotation for that reaction, and as a \emph{non-inferable gap} if no species in the pathway
does. Inferable gaps mark positions where cross-species evidence can prioritize candidate genes,
enzymes, or reactions (\texttt{wp:Conversion}) for further investigation, while non-inferable gaps
indicate pathway steps that currently require experimental characterization rather than annotation
transfer. To demonstrate this capability, we examined pathways representing distinct biological
scenarios, including species-specific specialized metabolism, medicinal alkaloid biosynthesis, and
specialized metabolites from non-model plants.

\note{Insert preprint Results lines 453--633 here: the four pathway case studies (capsaicin PWY-5710,
hyoscyamine/scopolamine PWY-7341, avenacin A-1 PWY-7476, monolignol PWY-361), the global annotation
matrix (29,061 combinations; 7,012 annotated, 13,922 inferable, 8,127 non-inferable) and the
\textit{A. thaliana} contribution analysis. Figures 3 and 4 belong in this subsection.}

\subsection*{Connecting plant metabolism to the Linked Open Data ecosystem through mapping and federation}

The PlantMetWiki SPARQL endpoint positions the resource as a key tool in the LOD ecosystem through
two complementary strategies: direct cross-resource identifier mapping (via BridgeDb), which
materializes links to external chemical databases inside the PlantMetWiki knowledge graph, and
federated queries, which reach any LOD resource at query time through shared identifiers, such as
InChIKeys\cite{leveau2019}. Together, these allow PlantMetWiki metabolites to be queried alongside
resources such as Wikidata, ChEBI\cite{malik2026}, and PubChem within a single SPARQL query. This
interoperability addresses a long-standing disconnect in plant biochemistry: pathway databases
capture how and in which species metabolites are produced, whereas chemical databases focus on
molecular properties, biological activities, toxicology, and applications. In Wikidata alone, the
production of atropine (Q26272) is not connected to the \textit{A. belladonna} plant (Q156091):
mapping and federation bridge these complementary resources.

\note{Insert preprint Results lines 652--860 here: BridgeDb coverage statistics and Figure 5, the
alpha-tomatine and theobromine case studies, and the Wikidata coverage / chemical novelty analysis
with Figure 6 (2,876 exact matches, 814 skeleton-only, 421 novel scaffolds).}

%==============================================================================
% DISCUSSION
%==============================================================================
\section*{Discussion}

\subsection*{PlantMetWiki enables FAIR multispecies representation of plant metabolism for cross-species comparison, pathway gap analysis, and candidate finding}

PlantMetWiki transforms species-centric pathway information into a FAIR multispecies knowledge graph
that exposes previously hidden relationships and annotation gaps across pathway, genomic, taxonomic,
and chemical resources. While Plant Metabolic Network remains the authoritative source of pathway
curation, PlantMetWiki provides a complementary semantic representation optimized for
interoperability, integration, and reuse, using standard RDF vocabularies, persistent identifiers,
machine-readable provenance, and SPARQL-based access. This LOD representation enables pathway
knowledge to be queried together with external biological, chemical, and genomic resources without
requiring duplication of those resources or manual data integration.

Semantic integration via existing vocabularies, such as the WikiPathways (\texttt{wp:}) one,
functions not only as a mechanism for interoperability but also as a framework for identifying
knowledge gaps across resources. Cross-species pathway queries revealed pathway positions lacking
annotation in individual species despite supporting evidence elsewhere in the graph. Integration with
biosynthetic gene cluster resources exposed clustered genomic regions that currently lack pathway
assignments. Federated querying against Wikidata revealed both representation-level discrepancies
arising from alternative stereochemical forms and 421 metabolites for which no structurally related
representation currently exists. These examples demonstrate that incompleteness can become visible
only when complementary knowledge sources are considered simultaneously. The distinction between
exact InChIKey matches and skeleton-level matches illustrates the importance of chemical
normalization when assessing interoperability between biological knowledge graphs, as many apparent
absences reflect alternative stereochemical or protonation representations rather than genuinely
missing chemistry. The reliable determination of stereochemistry and the subsequent reporting of
certainty and uncertainty of chiral centers in compounds are challenging aspects that LOD and
federation can help to normalize across resources with different curation depth. As stereochemistry
plays a key role in specialized compound bioactivity, with increasing reporting standards and
uptake\cite{cicek2024}, resources like PlantMetWiki can help to uncover specific stereochemical
patterns across multiple species.

Integration of biosynthetic gene cluster annotations illustrates how semantic representations can
connect traditionally separate views of plant metabolism. Pathway databases describe biochemical
transformations, whereas genome-mining resources focus on clustered genomic loci. By linking
PlantCyc-derived pathways to experimentally validated clusters from MIBiG and predicted clusters from
plantiSMASH, PlantMetWiki enables these complementary pathway representations and collections to be
explored jointly. The resulting crosslinks connect clustered and unclustered biosynthesis within a
common framework and provide a basis for investigating how genomic organization relates to metabolic
pathway structure. This demonstrates how FAIR knowledge graph approaches and Semantic Web
technologies can bridge resources that were originally developed for different biological questions
or with different search methods (such as genome mining in the case of plantiSMASH and MIBiG). Thanks
to these crosslinks, the PlantCyc pathways can now be connected to clustered regions in the genome,
which were not available either via PMN or MIBiG repositories. By linking clustered and unclustered
biosynthesis, PlantMetWiki will further help to understand the complex metabolic network in plants,
and their chromosomal arrangement.

A key advantage of PlantMetWiki is its ability to connect pathway knowledge with biosynthetic gene
cluster (BGC) information, bridging curated metabolic pathways with genome mining approaches, and
fill the knowledge gaps in complex and partially clustered biosynthesis. By integrating
experimentally validated clusters from MIBiG 4.0 and predicted clusters from plantiSMASH 2.0, the
resource enables users to explore relationships between clustered and unclustered biosynthetic routes
and to identify candidate genes associated with specific metabolic conversions and fills annotation
gaps across resources focusing on clustered or unclustered pathways. For instance, the alpha-tomatine
BGC has been key to characterize two related glycoalkaloid biosynthetic pathways involving
orthologous genes between tomato (\textit{Solanum lycopersicum}) and potato (\textit{Solanum
tuberosum}) and spanning across two biosynthetic loci: a cluster on chromosome 7 and an additional
duplicated region on chromosome 12\cite{itkin2013}. Because of this incomplete clustering,
alpha-tomatine was retired from the MIBiG 4.0 update, resulting in no available mapping of these loci
to plantiSMASH 2.0 output. The plantiSMASH clusterBLAST module reports the similarity of \textit{S.
lycopersicum} BGCs with ten other BGCs across species, to facilitate comparative analysis. By mapping
plantiSMASH-detected regions to PlantCyc pathway annotations, PlantMetWiki facilitates further
comparative studies involving pathways across species, resources, and genomic context.

PlantMetWiki provides a unified framework for integrating plant metabolic pathway knowledge across
diverse data types and species, enabling new opportunities for hypothesis generation in plant natural
product research across species from distributed biological knowledge. By combining curated metabolic
pathway annotations with Semantic Web technologies, PlantMetWiki facilitates the exploration of
biosynthetic pathways as interconnected, queryable graphs. This approach supports collaborative
discovery by allowing researchers to identify gaps in pathway knowledge, compare biosynthetic routes
across taxa, transfer knowledge from well-characterized species to less-studied taxa, and link
chemical occurrence data to underlying biosynthetic mechanisms. The resulting graph captures both
shared and species-specific aspects of metabolism, providing a reusable framework for studying
pathway evolution, annotation transfer, and specialized metabolism across the plant kingdom. In
addition, the availability of a SPARQL endpoint allows seamless integration with visualization and
analysis platforms such as Cytoscape\cite{shannon2003}, facilitating the exploration of metabolic and
regulatory networks in combination with user-provided omics data.

These findings demonstrate the power of interoperability for generating biologically meaningful
insights: by making pathway, taxonomic, genomic, and chemical information queryable within a common
framework, PlantMetWiki enables systematic assessment of annotation completeness across
complementary resources. Federated queries allow researchers to move seamlessly from biological
activity and toxicology to biosynthetic origin, genomic context, and species distribution, and
address completeness and curation of knowledge gaps.

\subsection*{Limitations: annotation quality of source databases, genome versioning, identifier mapping, and external endpoint availability}

Several limitations should be considered when using and interpreting the current resource. First,
PlantMetWiki inherits limitations from its source databases. Pathway completeness, species coverage,
and identifier consistency ultimately depend on the underlying curation available in PlantCyc, MIBiG,
plantiSMASH, and external reference resources. Second, cross-species inference to resolve annotation
gaps is limited to truly shared reaction steps across species, which is known to be only partial due
to the intrinsic fast-evolving nature of plant specialized metabolism\cite{weng2012,dadras2025}.
Third, integration across genomic resources remains constrained by differences in genome annotation
systems and identifier schemes: incompatible annotation frameworks can limit automated cross-resource
mapping\cite{monnahan2020}. For instance, for tomato (\textit{Solanum lycopersicum}), direct linking
of PlantCyc data with plantiSMASH and MIBiG annotations was not implemented because plantiSMASH uses
the NCBI RefSeq genome annotation (GCF\_036512215.1, generating LOC identifiers and gene symbols;
BioProject PRJDB16644), while PlantCyc uses the ITAG annotation\cite{sato2012} (generating Solyc
identifiers hosted in SolGenomics Network\cite{fernandez2015}). These two annotation systems define
different gene models for the same genome sequence and cannot be directly matched. BridgeDb can
support mapping annotation across resources, but the capacity of such tools needs to be constantly
updated to represent different sources and taxa, and it cannot bridge across truly different
annotations, such as gene annotations coming from different reference genomes, which is known as one
of the biggest challenges limiting the transferability of findings and integration of omics
experiments in the plant sciences field\cite{roeder2025}. Maintaining and extending current mapping
tools to additional species and genome annotations will be important for improving interoperability
across plant genomic resources. Additionally, the dependence on genome annotations introduces
challenges related to versioning and consistency, particularly in plant systems where genome
assemblies and annotations are frequently updated. Addressing this issue requires adherence to FAIR
data principles and explicit tracking of genome annotation versions, for example, through
standardized references to GenBank assemblies, to ensure reproducibility and transparency. The
PlantMetWiki knowledge graph provides a flexible framework to manage these evolving datasets but does
not eliminate the need for careful data curation.

Federated analyses depend on the continued availability and interoperability of external resources.
While LOD approaches reduce information duplication, they also rely on external services to maintain
stable identifiers and accessible query interfaces. This limitation is particularly apparent for
plantiSMASH and MIBiG, whose biological content is currently accessible primarily through web
interfaces rather than semantically modelled LOD, limiting the discoverability of the information
contained in these resources. As a result, PlantMetWiki can link users to biosynthetic gene clusters,
but cannot yet perform deeper federated queries across cluster-associated genes, metabolites,
taxonomic annotations, or chemical structures. The utility of federated analyses remains dependent on
the continued availability of public SPARQL endpoints such as Wikidata and the NCBI Taxonomy
project\cite{federhen2012}. These limitations underscore the importance of continued adoption of FAIR
principles, identifier harmonization, and LOD technologies across plant biology resources. Expanding
the ecosystem of queryable resources, for example, through public mirrors of databases such as ChEMBL
or the conversion of domain-specific datasets into LOD formats, will further enhance the utility of
Semantic Web platforms like PlantMetWiki.

\subsection*{Conclusions and Outlook: PlantMetWiki as a FAIR infrastructure for pathway knowledge}

PlantMetWiki provides a foundation for a continuously evolving, community-driven knowledge base for
plant biosynthesis. The integration of additional biological and omics layers will further enhance
the ability to model complex biosynthetic systems and generate testable hypotheses. The PlantMetWiki
framework can be extended to incorporate additional layers of biosynthetic regulation, including
transcription factor networks and other regulatory annotations, developmental stages, phenotypes,
gene homology, or orthologous relationships, as well as alternative genome mining or pathway-finding
tools beyond plantiSMASH. Experimental components such as tissues and developmental stages can be
integrated to advance the interpretation of omics and multi-omics layers at increased resolution, as
it has been demonstrated across paired omics studies\cite{wolters2024,barrios2026,jeon2020}. Such
integration highlights the potential of PlantMetWiki as a central hub for linking pathway knowledge
with emerging computational tools for biosynthetic discovery relying on multi-omics, such as
NPLinker\cite{geng2026} or MEANtools\cite{singhk2025}. Furthermore, PlantMetWiki has the potential to
connect to other knowledge graphs in the plant sciences, such as KnetMiner\cite{hassanipak2021}, for
gene and gene networks, AgroLD\cite{larmande2025}, for molecular and phenotypic information,
PlantConnectome\cite{lim2025}, for plant science literature, and Stress Knowledge
Map\cite{bleker2024}, for plant stress responses. Complementary to these resources, PlantMetWiki
shows the potential of Semantic Web technologies and LOD to transform various types of plant data
into actionable knowledge for hypothesis formulation and validation.

The modular architecture of PlantMetWiki enables iterative updates and continuous integration of new
data sources, including future releases of the source databases, Plant Metabolic Network, BridgeDb
identifier mappings, plantiSMASH predictions, and MIBiG annotations. Furthermore, the use of
federated queries enhances interoperability by allowing dynamic connections to external resources
that are regularly updated, such as Wikidata, ChEBI, PubMed, PubChem, and others. This approach can
help address common challenges in data integration, such as incomplete cross-referencing of
literature identifiers or missing metabolite annotations, by linking distributed knowledge across
databases without requiring centralized storage. Future work should focus on expanding species
coverage beyond current PlantCyc annotations and enabling joint querying of experimentally validated
and computationally predicted biosynthetic pathways.

The entire workflow underlying PlantMetWiki is transferable to other pathway collections from the
Pathway Tools ecosystem, such as MetaCyc and BioCyc organism databases: the conversion pipeline
operates on standard BioCyc flat-file exports and produces FAIR RDF representations with
machine-readable metadata and validation procedures. PlantMetWiki therefore represents both a
biological resource and a reusable methodology for transforming pathway databases into a fully
deployed RDF resource and interoperable knowledge graphs.

PlantMetWiki demonstrates how plant metabolic pathway knowledge can be represented as a FAIR
multispecies knowledge graph that supports integration across pathway, taxonomic, genomic, and
chemical domains. By combining semantic representation with interoperable identifiers and federated
querying, the resource enables comparative analyses, systematic assessment of annotation
completeness, and integration with complementary biological knowledge sources. Beyond the biological
content itself, PlantMetWiki provides a reproducible and transferable workflow for transforming
pathway databases into FAIR knowledge graphs, supporting broader reuse of pathway knowledge within
the life sciences LOD ecosystem.

%==============================================================================
% METHODS
%==============================================================================
\section*{Materials and Methods}

\subsection*{Generation of cross-species plant pathway representations}
\note{Insert preprint lines 1075--1128 (PlantCyc 17.0.0 retrieval, Cyc2Wiki extraction, grid-based
reaction layout, two-level taxonomy annotation, complex/monomer handling, pre-build validation,
citation transfer, GPML2021 XSD validation).}

\subsection*{RDF conversion, semantic enrichment, and BridgeDb mapping}
\note{Insert preprint lines 1133--1190 (gpml2rdf v4.0.4, PC/RC identifier assignment, parallelised
conversion, IRI scheme, BridgeDb build 20260102 and wp:bdb* predicates, taxonomy-extra and
properties-extra layers, three-level RDFlib validation, named graph structure).}

\subsection*{Extending pathway knowledge with crosslinks to biosynthetic gene clusters}
\note{Insert preprint lines 1196--1218 (BGC-first data model, pmw:BiosyntheticGeneCluster,
RO:0000051 has\_part, identifier normalisation to identifiers.org/tair.name/, SPARQL join strategy,
VoID metadata).}

\subsection*{Deployment, validation and accessibility}
\note{Insert preprint lines 1223--1281 (SNORQL UI and example query repository, QLever-based
federation, NCBITaxon MIREOT subsetting with ROBOT, VoID metadata design and discoverability at the
well-known VoID URI).}

%==============================================================================
% BACK MATTER
%==============================================================================
\section*{Data availability}

All relevant data are deposited and versioned in Zenodo at
\url{https://zenodo.org/communities/plantmetwiki/}, including relevant license and copyright
information. Tutorial pages are available at \url{https://pathway-lod.github.io/SPARQLTutorials/}.
Jupyter notebooks containing the steps to reproducibly generate the figures are available on GitHub
at \url{https://github.com/pathway-lod/gpml-to-rdf/tree/main/notebooks}. PlantMetWiki is listed on
bio.tools at \url{https://bio.tools/plantmetwiki}. Further project metadata is available on the
Nanodash platform (Knowledgepixels) at \url{https://w3id.org/spaces/plantmetwiki}.

\section*{Code availability}

All relevant code to this manuscript is deposited in GitHub at \url{https://github.com/pathway-lod}
and deposited in versioned releases in Zenodo at \url{https://zenodo.org/communities/plantmetwiki/}.

\section*{Supplementary tables}

\textbf{Supplementary Table S1.} Summary of taxa per NCBI ID represented in the GPML files.\\
\textbf{Supplementary Table S2.} Summary of taxonomic information captured by PlantMetWiki in the GPML representation.\\
\textbf{Supplementary Table S3.} Validation pipeline of GPML to RDF conversion.\\
\textbf{Supplementary Table S4.} Summary of extra properties captured by the third graph and not semantically modelled.\\
\textbf{Supplementary Table S5.} Summary of BGC mapping to PlantMetWiki from MIBiG and plantiSMASH.\\
\textbf{Supplementary Table S6.} Overview of the PlantMetWiki Virtuoso triplestore.\\
\textbf{Supplementary Table S7.} PlantMetWiki metabolites absent from Wikidata.

\section*{Author contributions}

Elena Del Pup: Conceptualization, Data curation, Formal analysis, Funding acquisition, Investigation,
Software, Writing -- original draft.
Max Muller: Data curation, Formal analysis, Software, Validation.
Egon L. Willighagen: Conceptualization, Formal analysis, Methodology, Software, Writing -- review \& editing.
Marvin Martens: Methodology, Software, Validation.
Marnix H. Medema: Funding acquisition, Project administration, Supervision, Writing -- review \& editing.
Denise Slenter: Conceptualization, Data curation, Formal analysis, Methodology, Software, Supervision,
Validation, Writing -- original draft, Writing -- review \& editing.
Justin J.J. van der Hooft: Conceptualization, Funding acquisition, Project administration, Supervision,
Writing -- review \& editing.

\note{The author list orders Martens before Willighagen; this paragraph orders Willighagen before
Martens. Confirm the intended order.}

\section*{Competing interests}

The authors declare the following financial interests and personal relationships which may be
considered as potential competing interests: J.J.J.vdH. is a member of the Scientific Advisory Board
of NAICONS Srl., Milano, Italy and consults for Corteva Agriscience, Indianapolis, IN, USA. M.H.M. is
a member of the Scientific Advisory Boards of Hexagon Bio and Hothouse Therapeutics Ltd. All other
authors declare to have no competing interests.

\section*{Acknowledgements}

We thank Sue Rhee and the Plant Metabolic Network curators for granting an open license that allows
reuse of their data, and for their valuable feedback. We thank Simon Shaw and Mitja Zdouc for
valuable advice on linking biosynthetic gene cluster information to PlantMetWiki. We also thank Arjan
Draisma and Harm Nijveen for advice on hosting the database and website for the SPARQL Explorer. We
thank Adriano Rutz for input on the use of Wikidata and LOTUS, and Leen Abraham for advice on
integrating plant reference genomes. Finally, we thank the Open Science and Open Data community for
building and maintaining the resources that PlantMetWiki can connect to for its queries.

\section*{Funding statement}

This work was supported by the FAIR data fund 4th edition from 4TU Research Data to E.D.P and by the
Netherlands Organization for Scientific Research (NWO) Vidi Grant VI.Vidi.213.183 to M.H.M.

\section*{AI usage statement}

During the preparation of this work, the authors used ChatGPT for brainstorming and obtaining
feedback on the manuscript; Claude and Claude Code for coding assistance and debugging; GitHub
Copilot to aid in code review; and Grammarly for copyediting. After using these tools, the authors
reviewed and edited the content as needed and take full responsibility for the content.

%==============================================================================
% FIGURES
%==============================================================================
\section*{Figures}

\begin{figure}[htbp]\centering
% \includegraphics[width=\textwidth]{figures/figure1.pdf}
\caption{\textbf{PlantMetWiki data transformation pipeline and architecture.}}
\label{fig:pipeline}\end{figure}

\begin{figure}[htbp]\centering
% \includegraphics[width=\textwidth]{figures/figure2.pdf}
\caption{\textbf{PlantMetWiki core graph model.}}
\label{fig:model}\end{figure}

\note{Add Figures 3--6 (capsaicin semantic model; metabolic annotation matrix and cross-species
inference potential; metabolite cross-reference coverage; Wikidata coverage and chemical novelty)
with their full preprint legends. Unlike Data Descriptors, Articles are not held to the
``no more than three figures'' guidance, so all six main figures can be retained.}

\bibliographystyle{unsrt}
\bibliography{references}

\end{document}
